\documentclass{article} %210 mm × 297 mm
\usepackage[utf8]{inputenc}
\usepackage[a4paper,left=1.5cm, right=1.5cm, top=2cm, bottom=2cm]{geometry}
\usepackage{graphicx}
%\usepackage{blindtext}
\usepackage{multirow}
\usepackage{mhchem}
\usepackage{url}
\usepackage{subfigure}
\usepackage{amsfonts,latexsym} % para tener disponibilidad de diversos simbolos
\usepackage{enumerate}
%\usepackage{booktabs}
\usepackage{float}
%\usepackage{threeparttable}
\usepackage{array,colortbl}
%\usepackage{ifpdf}
%\usepackage{rotating}
\usepackage{cite}
\usepackage{stfloats}
\usepackage{url}
\usepackage{listings}
\usepackage{fancyhdr}
\usepackage{biblatex}
\usepackage{color}
\addbibresource{sources.bib}
%\usepackage{hyperref}
%\usepackage{wrapfig}
\usepackage{array}
%\usepackage{multirow}
%\usepackage{adjustbox}
%\usepackage{nccmath}
\usepackage[dvipsnames]{xcolor}
\usepackage{tcolorbox}
\newtcolorbox{mybox}{colback=white,
colframe=black}
\newcommand{\titulo}{Abgabe 6; Diskrete Mathematik}
\newcommand{\nombre}{Liliana Sanfilippo}
\newcommand{\FCE}{\textbf{WiSe 2024/2025}}

 

\usepackage{fancyhdr}
\pagestyle{fancy}
\fancyhead{}
\fancyfoot{}
\fancyhead[LO]{\small\nombre}
\fancyhead[RE]{\small\titulo}
\fancyhead[LO]{\small\FCE}
\fancyfoot[RO,LE] {\thepage}
\usepackage[onehalfspacing]{setspace}

\setcounter{page}{1} %Número de la página, la editorial lo cambiará
\begin{document}
\begin{tabular}{p{12cm}}
{{\Large\bf\titulo}}\\
\vspace{0.2cm}\\
 {\large\nombre}\\
 \vspace{0.2cm}\\   
\end{tabular}
\section*{Präsenzaufgaben}
\subsection*{Aufgabe 1}
Es seien $\alpha = (13624)(587)(9)$ und $\beta = (15862)(394)(7)$. Geben Sie eine Permutation
$\sigma \in S9$ an, so dass $\sigma\alpha\sigma^{-1} = \beta$ gilt. Schreiben Sie $\sigma$ als Produkt disjunkter Zykeln. Wie
viele Möglichkeiten gibt es für die Wahl von $\sigma$?
\begin{mybox}
Ich fände es schön, wenn wir diese Aufgabe doch noch machen und besprechen könnten, denn ich verstehe sie nicht recht. 
\end{mybox}
\subsection*{Aufgabe 2}
Geben Sie alle Partitionen von n = 8 zusammen mit dem korrespondierenden Young
Diagramm an.  
\begin{mybox}
\begin{tabular}{cl}
      8 & xxxxxxxx \\ \hline
      7+1 & xxxxxxx \\
      & x\\ \hline
      6+2 & xxxxxx\\
      & xx\\\hline
      6+1+1 & xxxxxx\\
      & x\\
      & x\\\hline
      5+3 & xxxxx\\
      & xxx\\
\end{tabular}
\begin{tabular}{cl}
      5+2+1 & xxxxx \\
      & xx\\
      & x\\\hline
      5+1+1+1 & xxxxx \\
      & x \\
      & x \\
      & x \\\hline
      4+4 & xxxx\\
      & xxxx\\\hline
      4+3+1 & xxxx\\
      & xxx\\
      & x\\
\end{tabular}
\begin{tabular}{cl}
      4+2+2 & xxxx \\
      & xx\\
      & xx\\\hline
      4+2+1+1& xxxx\\
      & xx\\
      & x\\
      & x\\\hline
      4+1+1+1+1 & xxxx\\
      & x\\
      & x\\
      & x\\
      & x\\
\end{tabular}
\begin{tabular}{cl}
      3+3+2 & xxx \\
     & xxx\\
     & xx\\\hline
     3+3+1+1 & xxx\\
     & xxx\\
    &  x\\
     & x\\ \hline
     3+2+2+1 & xxx\\
     & xx\\
     & xx\\
     & x\\
\end{tabular}
\begin{tabular}{cl}
      3+2+1+1+1 & xxx\\
      & xx\\
      & x\\
      & x\\
      & x\\ \hline
      2+2+2+2 & xx\\
      & xx\\
      & xx\\
      & xx\\
\end{tabular}
\end{mybox}
\section*{Aufgaben zur Abgabe}
\subsection*{Aufgabe 1 (4 Punkte)}

Beweisen Sie, dass
\[
S(n, k) = \frac{1}{k!} 
\sum_{} 
\binom{n}{n_1, n_2, \ldots, n_k},
\]
wobei über alle Sequenzen \((n_1, \ldots, n_k)\) natürlicher Zahlen mit \(n_1, \ldots, n_k \geq 1\) und 
\(n_1 + n_2 + \cdots + n_k = n\) summiert wird.
\begin{mybox}
Wir analysieren die gegeben Formel und vergleichen anschließend die beiden Seiten. \\
$S(n,k)$ ist eine Stirling-Zahl zweiter Art und gibt die Anzahl der Möglichkeiten an, eine Menge mit $n$ Elementen in $k$ nichtleere, disjunkte Teilmengen zu zerlegen. \\
Multipliziert man beide Seiten der Gleichung mit $k!$, erhält man
\[
k!S(n, k) = 
\sum_{} 
\binom{n}{n_1, n_2, \ldots, n_k}
\]
Wir wissen bereits, dass $k!S(n, k)$ dia Anzahl der surjektiven Abbildungen zwischen einer Menge von $n$ Objekten auf eine Menge von $k$ Fächer ist und kennen $\binom{n}{n_1, n_2, \ldots, n_k}$ als Multinomialkoeffizienten. \\ \newline
Aus dem Skript wissen wir, dass der Multinomialkoeffizient angibt, wie viele Möglichkeiten es gibt eine Menge mit $n$ Elementen in $k$ geordnete Teilmengen aufzuteilen, bei denen jeweils $n_x$ Elemente auf eine Zahl $x$ abgebildet werden. Dabei betrachten wir die Anzahl mit Reihenfolge. \\
Summiert man alle möglichen Multinomialkoeffizienten, erhält man die Anzahl aller möglichen Zerlegungen von $n$ Elementen in $k$ Teilmengen unter der Berücksichtigung der Reihenfolgen der Elemente und das entspricht $k!S(n, k)$. 
Damit ist folgendes bewiesen: 
\[
k!S(n, k) = 
\sum_{} 
\binom{n}{n_1, n_2, \ldots, n_k}
\]
Und es folgt die Korrektheit der Formel 
\[
S(n, k) = \frac{1}{k!} 
\sum_{} 
\binom{n}{n_1, n_2, \ldots, n_k}
\]


\end{mybox}
\subsection*{Aufgabe 2 (4 Punkte)}

Zeigen Sie, dass für alle natürlichen Zahlen \(a, b, c\) mit \(a + b + c = n\) gilt:
\[
\binom{n}{a, b, c} = 
\binom{n-1}{a-1, b, c} + 
\binom{n-1}{a, b-1, c} + 
\binom{n-1}{a, b, c-1}.
\]

Geben Sie eine analoge Formel für allgemeine Multinomialzahlen an.

\begin{mybox}
$\binom{n}{a, b, c}$ ist ein Multinomialkoeffizient. Dieser spezielle Multinomialkoeffizient gibt die Anzahl der Möglichkeiten an $n$ Objekte auf drei Teilmengen mit $a$, $b$ und $c$ vielen Elementen abzubilden. \\
$\binom{n-1}{a-1, b, c} + 
\binom{n-1}{a, b-1, c} + 
\binom{n-1}{a, b, c-1}$ ist eine Rekursionsformel und erlaubt es, dass man zuerst nur das $n$-te Objekt und dessen Zuweisung betrachtet. Dabei gibt es drei Möglichkeiten. 
\begin{enumerate}
    \item Das $n$-te Objekt wurde der Teilmenge mit $a$ vielen Elementen zugeordnet \\
    $\rightarrow$ Dann werden die restlichen $n-1$ Objekte auf drei Teilmengen mit $a-1$, $b$ und $c$ vielen Elementen abgebildet; Es gibt $\binom{n-1}{a-1, b, c}$ viele Möglichkeiten
    \item Das $n$-te Objekt wurde der Teilmenge mit $b$ vielen Elementen zugeordnet\\
    $\rightarrow$ Dann werden die restlichen $n-1$ Objekte auf drei Teilmengen mit $a$, $b-1$ und $c$ vielen Elementen abgebildet; Es gibt $\binom{n-1}{a, b-1, c}$ viele Möglichkeiten
    \item Das $n$-te Objekt wurde der Teilmenge mit $c$ vielen Elementen zugeordnet\\
    $\rightarrow$ Dann werden die restlichen $n-1$ Objekte auf drei Teilmengen mit $a$, $b$ und $c-1$ vielen Elementen abgebildet; Es gibt $\binom{n-1}{a, b, c-1}$ viele Möglichkeiten
\end{enumerate}
Addiert man die drei sich daraus ergebenden Anzahlen von Möglichkeiten, erhält man die Gesamtanzahl von Möglichkeiten $n$ Objekte auf drei Teilmengen mit $a$, $b$ und $c$ vielen Elementen abzubilden.\\ \newline
Diese Summierung der Anzahlen der Partitionen abhängig davon wo das $n$-te Element ist, kann man verallgemerinern als 
\[
\binom{n}{r_1, r_2, \dots, r_k} = \sum_{i=1}^{k} \binom{n-1}{r_1, \dots, r_{i-1}, r_i-1, r_{i+1}, \dots, r_k} 
\]
\end{mybox}

\subsection*{Aufgabe 3 (4 Punkte)}

Sei \(n = a_1 + \cdots + a_k\) eine Partition. Sei \(b_i := |\{j \mid a_j \geq i\}|\). Die konjugierte Partition ist 
\(b_1 + b_2 + \cdots\).

\begin{enumerate}
    \item[(a)] Schreiben Sie alle Partitionen von \(n = 6\) und \(n = 7\) auf. Für jede Partition von \(n = 6\) und \(n = 7\), geben Sie die konjugierte Partition an.
   \begin{mybox}
    \begin{tabular}{cc|c}
        n = 6 & Partitionen & konjugierte Partitionen\\
         & 6 & 1+1+1+1+1+1 \\
         & 5+1 & 2+1+1+1+1 \\
         & 4+2 & 2+2+1+1 \\
         & 4+1+1 & 3+1+1+1 \\
         & 3+3 & 2+2+2\\
         & 3+2+1  & 3+2+1\\
         & 2 +2+2 & 3+3 \\
         & 2+2+1+1&  4+2\\
         & 1+1+1+1+1+1 & 6\\
    \end{tabular}
    \end{mybox}
    \begin{mybox}
    \begin{tabular}{cc|c}
    n = 7 & Partitionen & konjugierte Partitionen\\
         & 7 &   1+1+1+1+1+1+1\\
         & 6+1& 2+1+1+1+1+1+1\\
         & 5+2& 2+2+1+1+1 \\
         & 5+1+1 & 3+1+1+1+1\\
         &4+3& 2+2+2+1 \\
         &4+2+1 & 3+2+1+1 \\
         & 4+1+1+1& 4+1+1+1 \\
         &3+2+1+1& 4+2+1\\
         & 3+1+1+1+1& 5+1+1 \\
         &2+2+2+1&4+3\\
         &2+2+1+1+1& 5+2 \\
         &2+1+1+1+1+1& 6+1\\
         & 1+1+1+1+1+1+1& 7\\
    \end{tabular} 
   \end{mybox}
    \item[(b)] Zeigen Sie, dass \(n = b_1 + b_2 + \cdots\), das heißt, \(b_1 + b_2 + \cdots\) ist ebenfalls eine Partition von \(n\).
    \begin{mybox}
    Wir wollen zeigen, dass die Summe aller $b_i$ eine Partition von $n$ ist. \\
    $b_1 + b_2 + \cdots$ beschreibt wie viele Teile der Partition $a_1 + a_2 + \cdots + a_k$ größer oder gleich $i$ sind. Für jedes $i$ wird also jedes Element $a_j$ genau einmal in dem $b_i$ betrachtet und (wenn es $\geq i$ ist) dazu gezählt. Dadurch wird jedes $a_j$ genau so oft in die Summe der $b_i$s einbezogen, wie der Wert von $a_j$ in der Partition ist. Und dadurch ist die Summe identisch zu $n$ und somit eine Partition von $n$.
    \end{mybox}
\end{enumerate}

\subsection*{Aufgabe 4 (4 Punkte)}

Sei \(p_k(n)\) die Anzahl der Partitionen von \(n\) mit genau \(k\) Summanden. Beweisen Sie die folgende Rekursionsformel:
\[
p_k(n) = p_k(n - k) + p_{k-1}(n - k) + \cdots + p_1(n - k), \quad n > k \geq 1.
\]
\begin{mybox}
Wie bei Aufgabe 2 müssen wir uns anschauen welche Fälle es bei der Zerlegung von $p_k(n)$ in eine Rekursionsformel geben kann. 
\begin{enumerate}
    \item Der größte Summand einer Partition ist gleich $k$ \\
    $\rightarrow$ Die restlichen $k-1$ Summanden bilden dann eine Partition von $n-k$. Die Anzahl dieser Partitionen ist $p_{k-1}(n-k)$
    \item Der größte Summand der Partition ist kleiner als $k$.  \\
    $\rightarrow$ Dann sind die Partitionen von $n$ mit $k$ Summanden identisch mit den Partitionen von $n-k$ in maximal $k-1$ Summanden, da der größte Summand kleiner als $k$ ist. \\
    Die Anzahl der Partitionen in Summanden kleiner als $k$ ist $\sum_{i=1}^{k-1} p_{i}(n-k)$, also $p_{k-1}(n-k) + p_{k-2}(n-k) \cdots + p_{1}(n-k)$. 
\end{enumerate} 
Nimmt man all diese Anzahlen zusammen (und deckt somit alle Fälle ab), erhält man die Gesamtanzahl der Partitionen von $n$ mit genau $k$ Summanden. 
\[
p_{k}(n-k) + \sum_{i=1}^{k-1} p_{i}(n-k) = p_{k}(n-k) + p_{k-1-i}(n-k) + p_{k-2}(n-k) \cdots + p_{1}(n-k)
\]
\end{mybox}
\end{document}

  %Wenn der größte Summand kleiner als $k$ ist, dann muss der größte Summand mindestens 1 sein und 